\section{结论}

本文对ResNet论文《Deep Residual Learning for Image Recognition》进行了全面解析。文章首先指出了深度卷积网络面临的退化问题，即随着层数增加训练误差上升的优化困境。为攻克这一难题，He等人提出了残差学习框架，利用恒等快捷连接将堆叠层需要学习的映射转化为残差映射，使得网络能够轻松逼近恒等变换。

实验表明，残差结构不仅消除了退化现象，还支持训练极深的网络（如152层、1202层），并在ImageNet、CIFAR-10等多个基准测试中取得了创纪录的成绩。进一步的消融研究验证了恒等连接优于投影连接，且残差块设计简洁、参数高效。

从理论角度看，残差网络通过梯度中的直通路保证了反向传播的畅通，同时可被解释为多个浅层网络的隐式集成，并且前向时更容易学习微小扰动。这些特性共同解释了残差学习的成功。

最后，ResNet深刻影响了后续深度学习架构的发展，催生了ResNeXt、DenseNet、Wide ResNet等经典变体，其残差连接已成为现代神经网络的标准组件。正如论文所言：“深度是好的，但残差让深度成为可能。” 残差学习不仅是一项技术突破，更是一种设计思维：面对深层网络优化难题，简单而直接的恒等映射即可化繁为简，释放深度的力量。